\documentclass[11pt]{article}
\usepackage{mathblatt}
\def\mitzone{1}
\begin{document}
\blattkopf{Lineare Funktionen}{Gesamt}
\weit
\verzeichniszeile{\verz{zone}{Kennst du schon (Nr.~1–8)} \verztrenn \verz{e1}{1 Proportionale Funktion (Nr.~9–11)} \verztrenn \verz{e2}{2 Lineare Funktion f(x) = m·x + n (Nr.~12–18)} \verztrenn \verz{e3}{3 Punkte und Werte (Nr.~19–23)} \verztrenn \verz{e4}{4 Gleichung bestimmen (Nr.~24–27)} \verztrenn \verz{e5}{5 Anwendungen (Nr.~28–31)} \verztrenn \verz{abhaken}{Das kann ich}}
% Zone „Das kennst du schon“ – aus bank/lineare-funktionen/zone.jsonl (zusammenbau v0.1)
\einheitenkopf[zone][Kennst du schon]{Das kennst du schon}
\setcounter{aufgabe}{0}

%% TODO Titel: Ich-kann-Satz fehlt in der Bank (Platzhalter: Kettenname) – Nr. 1
%% TODO Anweisung: ein Satz über der ersten Teilaufgabe fehlt in der Bank – Nr. 1
\begin{aufgabe}{Koordinaten lesen und eintragen, 4 Quadranten, (x|y)-Reihenfolge}
\begin{teile}
\teil Punkt A – welche Koordinaten? Lies ab.

\begin{ksys}[xmin=-5,xmax=5,ymin=-5,ymax=5,ablesen] \punkt{4}{3}{A} \end{ksys}
A( \leerfeld | \leerfeld )
\teil Punkt C – welche Koordinaten? Lies ab.

\begin{ksys}[xmin=-5,xmax=5,ymin=-5,ymax=5,ablesen] \punkt{-3}{-4}{C} \end{ksys}
C( \leerfeld | \leerfeld )
\end{teile}
\end{aufgabe}

%% TODO Titel: Ich-kann-Satz fehlt in der Bank (Platzhalter: Kettenname) – Nr. 2
%% TODO Anweisung: ein Satz über der ersten Teilaufgabe fehlt in der Bank – Nr. 2
\begin{aufgabe}{Proportionale Zuordnung erkennen und hochrechnen (Dreisatz)}
\begin{teile}
\teil Preis für 5 kg? 2 kg Äpfel kosten 4 €. Wie viel kosten 5 kg? \leerfeld[€]
\teil Preis für 3 Liter? 4 Liter Saft kosten 6 €. Wie viel kosten 3 Liter? \leerfeld[€]
\end{teile}
\end{aufgabe}

%% TODO Titel: Ich-kann-Satz fehlt in der Bank (Platzhalter: Kettenname) – Nr. 3
%% TODO Anweisung: ein Satz über der ersten Teilaufgabe fehlt in der Bank – Nr. 3
\begin{aufgabe}{Negative Zahlen multiplizieren und dividieren}
\begin{teile}
\teil Rechne: $(-3) \cdot 4$
\teil Rechne: $(-1{,}5) \cdot (-4)$
\end{teile}
\end{aufgabe}

%% TODO Titel: Ich-kann-Satz fehlt in der Bank (Platzhalter: Kettenname) – Nr. 4
%% TODO Anweisung: ein Satz über der ersten Teilaufgabe fehlt in der Bank – Nr. 4
\begin{aufgabe}{Brüche als Steigung (½, −¾), Bruch mal ganze Zahl}
\begin{teile}
\teil Rechne: $\frac{1}{2} \cdot 8$
\teil Rechne: $-\frac{3}{4} \cdot 8$
\end{teile}
\end{aufgabe}

%% TODO Titel: Ich-kann-Satz fehlt in der Bank (Platzhalter: Kettenname) – Nr. 5
%% TODO Anweisung: ein Satz über der ersten Teilaufgabe fehlt in der Bank – Nr. 5
\begin{aufgabe}{Lineare Gleichung zweischrittig lösen (null gleich einem zweischrittigen Term)}
\begin{gleichungsraster}[2]
\gl{\text{Löse: $0 = 2x - 8$}} &
\gl{\text{Löse: $0 = -4x + 12$}} \\
\end{gleichungsraster}
\end{aufgabe}

%% TODO Titel: Ich-kann-Satz fehlt in der Bank (Platzhalter: Kettenname) – Nr. 6
%% TODO Anweisung: ein Satz über der ersten Teilaufgabe fehlt in der Bank – Nr. 6
\begin{aufgabe}{Terme zusammenfassen}
\begin{teile}
\teil Fasse zusammen: $3x + 4x$
\teil Fasse zusammen: $2x + 5 - 4x + 1$
\end{teile}
\end{aufgabe}

%% TODO Titel: Ich-kann-Satz fehlt in der Bank (Platzhalter: Kettenname) – Nr. 7
%% TODO Anweisung: ein Satz über der ersten Teilaufgabe fehlt in der Bank – Nr. 7
\begin{aufgabe}{Negative Zahlen multiplizieren und dividieren – Fehler finden}
\begin{teile}
\teil Ich finde den Fehler: Mia rechnet $-3 \cdot (-5) + 2 = -13$. Benenne den Fehler und rechne richtig.
\end{teile}
\end{aufgabe}

%% TODO Titel: Ich-kann-Satz fehlt in der Bank (Platzhalter: Kettenname) – Nr. 8
%% TODO Anweisung: ein Satz über der ersten Teilaufgabe fehlt in der Bank – Nr. 8
\begin{aufgabe}{Negative Zahlen multiplizieren und dividieren – selbst rechnen}
\begin{teile}
\teil Rechne: $-6 \cdot (-4) + 5$
\end{teile}
\end{aufgabe}

\clearpage
% Einheit 1 von 5 – Proportionale Funktion (bank/lineare-funktionen/e1.jsonl, zusammenbau v0.1)
%% TODO Zweigzeile Teil 1: was hier gelernt wird (Satz in Schülersprache aus dem Einheitstitel) – Einheit 1
\einheitenkopf[e1]{Einheit 1 von 5 · Proportionale Funktion}
\zweigzeile{ab Kl. 8 · P10 oft}
\setcounter{aufgabe}{8}

%% TODO Titel: Ich-kann-Satz fehlt in der Bank (Platzhalter: Kettenname) – Nr. 9
%% TODO Anweisung: ein Satz über der ersten Teilaufgabe fehlt in der Bank – Nr. 9
\begin{aufgabe}{Proportionale Funktion}
%% TODO \rechenplatz{n}: Schrittzahl des Grundfalls fehlt in der Bank (nur bei fester Schreibform, 2.3 b)
\begin{teile}
\teil Welcher Faktor? Mit welcher Zahl wird jeder x-Wert multipliziert?

\wertetabelle[3,6,9,12]{x}{y}{1,2,3,4}
Faktor: \leerfeld
\teil Wertepaare und Gerade: Berechne die y-Werte zu $y = 3x$ für $x = 0, 1, 2$ und zeichne die Ursprungsgerade.

\begin{ksys}[xmin=-1,xmax=3,ymin=-1,ymax=7] \end{ksys}
\teil y-Wert bei $x = 3$? Lies am Graphen von $y = 2x$ ab und prüfe mit der Rechnung.

\begin{ksys}[xmin=-1,xmax=5,ymin=-1,ymax=8,ablesen] \gerade{2}{0}{g} \end{ksys}
y = \leerfeld
\teil Wertepaare und Gerade: Berechne die y-Werte zu $y = 1{,}5x$ für $x = 0, 2, 4$ und zeichne die Ursprungsgerade.

\begin{ksys}[xmin=-1,xmax=5,ymin=-1,ymax=7] \end{ksys}
\teil Wertepaare und Gerade: Berechne die y-Werte zu $y = \frac{1}{3}x$ für $x = 0, 3, 6$ und zeichne die Ursprungsgerade.

\begin{ksys}[xmin=-1,xmax=7,ymin=-1,ymax=3] \end{ksys}
\teil Wertepaare und Gerade: Berechne die y-Werte zu $y = -3x$ für $x = 0, 1, 2$ und zeichne die Ursprungsgerade.

\begin{ksys}[xmin=-1,xmax=3,ymin=-7,ymax=1] \end{ksys}
\teil Proportional? Entscheide mit den Quotienten $y : x$. Wenn ja: Gib die Gleichung an.

\wertetabelle[7,14,21,28]{x}{y}{1,2,3,4}
\janein \quad y = \leerfeld
\teil Ich finde den Fehler: Ben bestimmt den Faktor der Tabelle so: $10 - 2 = 8$, also $y = 8x$. Benenne den Fehler und gib die richtige Gleichung an.

\wertetabelle[10,20,30]{x}{y}{2,4,6}
\teil Wasser im Becken: In ein leeres Becken laufen je Minute 12 Liter Wasser. Stelle eine Gleichung für die Wassermenge y (in Litern) nach x Minuten auf, zeichne den Graphen für 0 bis 5 Minuten und erkläre, was der Faktor im Sachzusammenhang bedeutet.

\begin{ksys}[xmin=0,xmax=6,ymin=0,ymax=80,xstep=1,ystep=10,xlabel=x in min,ylabel=y in l] \end{ksys}
y = \leerfeld
\end{teile}
\end{aufgabe}

%% TODO Titel: Ich-kann-Satz fehlt in der Bank (Platzhalter: Kettenname) – Nr. 10
%% TODO Anweisung: ein Satz über der ersten Teilaufgabe fehlt in der Bank – Nr. 10
\begin{aufgabe}{Dreisatz als Kontrolle}
\begin{teile}
\teil Preis für 6 m? 4 m Kabel kosten 10 €. Ole rechnet mit $y = 2{,}5x$ (y Preis in €, x Länge in m). Kontrolliere mit dem Dreisatz den Preis für 6 m. \leerfeld[€]
\end{teile}
\end{aufgabe}

%% TODO Titel: Ich-kann-Satz fehlt in der Bank (Platzhalter: Kettenname) – Nr. 11
%% TODO Anweisung: ein Satz über der ersten Teilaufgabe fehlt in der Bank – Nr. 11
\begin{aufgabe}{Proportionale Funktion – Fehler finden, Begründen, Darstellungswechsel, Anwendung}
\begin{teile}
\teil Ich finde den Fehler: Zu $y = 3x$ trägt Jonas für $x = 2$ den Punkt $(6|2)$ ein. Benenne den Fehler und gib den richtigen Punkt an.
\teil Proportional? Eine Gerade schneidet die y-Achse oberhalb des Ursprungs. Kann sie zu einer proportionalen Funktion gehören? Begründe ohne Rechnung.
\teil Gleichung zum Graphen? Die Ursprungsgerade g geht durch den Punkt P. Gib ihre Gleichung an.

\begin{ksys}[xmin=-1,xmax=5,ymin=-1,ymax=7,ablesen] \gerade{1.5}{0}{g} \punkt{2}{3}{P} \end{ksys}
y = \leerfeld
\teil Reicht das Geld? Ein Liter Diesel kostet 1{,}70 €, der Preis y (in €) für x Liter ist $y = 1{,}7x$. Lisa hat 60 € dabei. Reicht das für 35 Liter? \janein
\end{teile}
\end{aufgabe}

\clearpage
% Einheit 2 von 5 – Lineare Funktion f(x) = m·x + n (bank/lineare-funktionen/e2.jsonl, zusammenbau v0.1)
%% TODO Zweigzeile Teil 1: was hier gelernt wird (Satz in Schülersprache aus dem Einheitstitel) – Einheit 2
\einheitenkopf[e2]{Einheit 2 von 5 · Lineare Funktion f(x) = m·x + n}
\zweigzeile{ab Kl. 8 · P10 oft}
\setcounter{aufgabe}{11}

%% TODO Titel: Ich-kann-Satz fehlt in der Bank (Platzhalter: Kettenname) – Nr. 12
%% TODO Anweisung: ein Satz über der ersten Teilaufgabe fehlt in der Bank – Nr. 12
\begin{aufgabe}{m und n in der Gleichung markieren}
\begin{teile}
\teil m und n? Unterstreiche m, kreise n ein: $f(x) = 3x - 5$ m = \leerfeld , n = \leerfeld
\end{teile}
\end{aufgabe}

%% TODO Titel: Ich-kann-Satz fehlt in der Bank (Platzhalter: Kettenname) – Nr. 13
%% TODO Anweisung: ein Satz über der ersten Teilaufgabe fehlt in der Bank – Nr. 13
\begin{aufgabe}{Steigungsrichtung ohne Zeichnen}
\begin{teile}
\teil Steigt oder fällt die Gerade $f(x) = 3x - 4$? Kreuze an. \\ \kreuz{steigt} \\ \kreuz{fällt}
\end{teile}
\end{aufgabe}

%% TODO Titel: Ich-kann-Satz fehlt in der Bank (Platzhalter: Kettenname) – Nr. 14
%% TODO Anweisung: ein Satz über der ersten Teilaufgabe fehlt in der Bank – Nr. 14
\begin{aufgabe}{Steigungsdreieck lesen}
\begin{teile}
\teil Einen Schritt nach rechts – wie viel nach oben oder unten? Lies am Steigungsdreieck ab.

\begin{ksys}[xmin=-4,xmax=4,ymin=-5,ymax=5,ablesen] \gerade{3}{-1}{g} \steigungsdreieck{0}{-1}{3} \end{ksys}
\leerfeld nach \leerfeld \leerfeld
\end{teile}
\end{aufgabe}

%% TODO \verfahren{…}: Verfahrensüberschrift in Sachsprache fehlt in der Bank (2.3 b)
%% TODO Titel: Ich-kann-Satz fehlt in der Bank (Platzhalter: Kettenname) – Nr. 15
%% TODO Anweisung: ein Satz über der ersten Teilaufgabe fehlt in der Bank – Nr. 15
\begin{aufgabe}{Graph zeichnen}
%% TODO \rechenplatz{n}: Schrittzahl des Grundfalls fehlt in der Bank (nur bei fester Schreibform, 2.3 b)
\begin{teile}
\teil Wertetabelle: Fülle die Tabelle zu $f(x) = 3x - 2$ aus.

\wertetabelle{x}{f(x)}{-1,0,1,2}
\teil Punkte eintragen: Trage die Wertepaare $(-2|-2)$, $(-1|1)$, $(0|4)$, $(1|7)$ der Funktion $f(x) = 3x + 4$ ins Koordinatensystem ein.

\begin{ksys}[xmin=-4,xmax=4,ymin=-4,ymax=8] \end{ksys}
\teil Gerade durch Punkte: Berechne zwei Punkte von $f(x) = 4x + 2$, trage sie ein und zeichne die Gerade durch sie.

\begin{ksys}[xmin=-4,xmax=4,ymin=-4,ymax=7] \end{ksys}
\teil Mit n und Steigungsdreieck: Zeichne die Gerade zu $f(x) = 3x + 5$: Markiere n auf der y-Achse, gehe 1 nach rechts und m nach oben.

\begin{ksys}[xmin=-4,xmax=4,ymin=-4,ymax=9] \end{ksys}
\teil Mit n und Steigungsdreieck: Zeichne die Gerade zu $f(x) = -3x + 4$.

\begin{ksys}[xmin=-4,xmax=4,ymin=-4,ymax=5] \end{ksys}
\teil Mit n und Steigungsdreieck: Zeichne die Gerade zu $f(x) = \frac{1}{2}x + 3$.

\begin{ksys}[xmin=-4,xmax=4,ymin=-4,ymax=5] \end{ksys}
\teil Mit n und Steigungsdreieck: Zeichne die Gerade zu $f(x) = 2x - 4$.

\begin{ksys}[xmin=-4,xmax=4,ymin=-5,ymax=4] \end{ksys}
\teil Graph zeichnen: Zeichne den Graphen der Funktion f mit $f(x) = -4x + 3$ in das Koordinatensystem. \hfill (P10 2026 FOR)

\begin{ksys}[xmin=-4,xmax=4,ymin=-4,ymax=4] \end{ksys}
\end{teile}
\end{aufgabe}

%% TODO \verfahren{…}: Verfahrensüberschrift in Sachsprache fehlt in der Bank (2.3 b)
%% TODO Titel: Ich-kann-Satz fehlt in der Bank (Platzhalter: Kettenname) – Nr. 16
%% TODO Anweisung: ein Satz über der ersten Teilaufgabe fehlt in der Bank – Nr. 16
\begin{aufgabe}{Ablesen}
%% TODO \rechenplatz{n}: Schrittzahl des Grundfalls fehlt in der Bank (nur bei fester Schreibform, 2.3 b)
\begin{teile}
\teil n? Lies den y-Achsenabschnitt n der Geraden g ab.

\begin{ksys}[xmin=-5,xmax=5,ymin=-5,ymax=5,ablesen] \gerade{1}{3}{g} \end{ksys}
n = \leerfeld
\teil m? Lies die Steigung m der Geraden g mit einem Steigungsdreieck ab.

\begin{ksys}[xmin=-5,xmax=5,ymin=-5,ymax=5,ablesen] \gerade{2}{-3}{g} \end{ksys}
m = \leerfeld
\teil m? Lies die Steigung m der Geraden g mit einem Steigungsdreieck ab.

\begin{ksys}[xmin=-5,xmax=5,ymin=-5,ymax=5,ablesen] \gerade{-2}{1}{g} \end{ksys}
m = \leerfeld
\teil m? Lies die Steigung m der Geraden g ab. Gehe so weit nach rechts, bis du wieder auf einer Gitterlinie bist.

\begin{ksys}[xmin=-5,xmax=5,ymin=-5,ymax=5,ablesen] \gerade{0.5}{1}{g} \end{ksys}
m = \leerfeld
\teil Welche Gleichung gehört zur Geraden? Kreuze an. \\ \kreuz{$f(x) = 3x - 4$} \\ \kreuz{$f(x) = -4x + 3$} \\ \kreuz{$f(x) = 3x + 4$} \\ \kreuz{$f(x) = -3x - 4$}

\begin{ksys}[xmin=-5,xmax=5,ymin=-5,ymax=5,ablesen] \gerade{3}{-4}{g} \end{ksys}
\teil Welche Gleichung gehört zu welcher Geraden? Ordne den Geraden f, g und h je eine der Gleichungen zu: $y = -3x + 3$; $y = -\frac{1}{3}x + 3$; $y = -3$; $y = 2x + 3$; $y = -3x$; $y = 2x - 3$. \hfill (P10 2019 OS) \\

\begin{ksys}[xmin=-5,xmax=5,ymin=-5,ymax=5,ablesen] \gerade{-3}{3}{f} \gerade{2}{3}{g} \gerade{0}{-3}{h} \end{ksys}
f: \leerfeld \quad g: \leerfeld \quad h: \leerfeld
\end{teile}
\end{aufgabe}

%% TODO Titel: Ich-kann-Satz fehlt in der Bank (Platzhalter: Kettenname) – Nr. 17
%% TODO Anweisung: ein Satz über der ersten Teilaufgabe fehlt in der Bank – Nr. 17
\begin{aufgabe}{Parameter deuten (steigend, fallend, parallel, Sonderfall m = 0)}
\begin{teile}
\teil Welche zwei Geraden sind parallel? $f(x) = 4x - 1$, $g(x) = -4x + 5$, $h(x) = 4x + 6$
\end{teile}
\end{aufgabe}

%% TODO Titel: Ich-kann-Satz fehlt in der Bank (Platzhalter: Kettenname) – Nr. 18
%% TODO Anweisung: ein Satz über der ersten Teilaufgabe fehlt in der Bank – Nr. 18
\begin{aufgabe}{Graph zeichnen – Fehler finden, Begründen, Darstellungswechsel}
\begin{teile}
\teil Ich finde den Fehler: Paul liest am Steigungsdreieck ab: 1 nach rechts, 3 nach oben, also $m = \frac{1}{3}$. Benenne den Fehler und gib m richtig an.

\begin{ksys}[xmin=-4,xmax=4,ymin=-5,ymax=5,ablesen] \gerade{3}{-2}{g} \steigungsdreieck{0}{-2}{3} \end{ksys}
\teil Welche Gerade ist steiler: $f(x) = 3x - 7$ oder $g(x) = 5x + 4$? Begründe ohne Rechnung.
\teil Gleichung zur Geraden? Gib die Gleichung der Geraden g an.

\begin{ksys}[xmin=-5,xmax=5,ymin=-5,ymax=5,ablesen] \gerade{-1}{3}{g} \end{ksys}
f(x) = \leerfeld
\end{teile}
\end{aufgabe}

\clearpage
% Einheit 3 von 5 – Punkte und Werte (bank/lineare-funktionen/e3.jsonl, zusammenbau v0.1)
%% TODO Zweigzeile Teil 1: was hier gelernt wird (Satz in Schülersprache aus dem Einheitstitel) – Einheit 3
\einheitenkopf[e3]{Einheit 3 von 5 · Punkte und Werte}
\zweigzeile{ab Kl. 8 · P10 oft}
\setcounter{aufgabe}{18}

%% TODO Titel: Ich-kann-Satz fehlt in der Bank (Platzhalter: Kettenname) – Nr. 19
%% TODO Anweisung: ein Satz über der ersten Teilaufgabe fehlt in der Bank – Nr. 19
\begin{aufgabe}{Punkt in Gleichung einsetzen}
\begin{teile}
\teil Einsetzen: Setze $x = 4$ in $3x - 2$ ein und rechne aus.
\end{teile}
\end{aufgabe}

%% TODO Titel: Ich-kann-Satz fehlt in der Bank (Platzhalter: Kettenname) – Nr. 20
%% TODO Anweisung: ein Satz über der ersten Teilaufgabe fehlt in der Bank – Nr. 20
\begin{aufgabe}{Funktionswert}
%% TODO \rechenplatz{n}: Schrittzahl des Grundfalls fehlt in der Bank (nur bei fester Schreibform, 2.3 b)
\begin{teile}
\teil $f(4)$? Berechne den Funktionswert von $f(x) = 2x + 5$ an der Stelle $x = 4$. f(4) = \leerfeld
\teil $f(-3)$? Berechne den Funktionswert von $f(x) = 4x - 1$ an der Stelle $x = -3$. f(-3) = \leerfeld
\teil $f(1{,}5)$? Berechne den Funktionswert von $f(x) = 4x - 3$ an der Stelle $x = 1{,}5$. f(1{,}5) = \leerfeld
\end{teile}
\begin{gleichungsraster}[2]
\gl{\text{Welches x? Für welches x gilt $f(x) = 11$ bei $f(x) = 2x + 3$?}} &
\end{gleichungsraster}
\begin{teile}
\teil Liegt der Punkt? Prüfe, ob $P(3|8)$ auf der Geraden $f(x) = 3x - 1$ liegt. \janein
\end{teile}
\begin{gleichungsraster}[2]
\gl{\text{Nullstelle? Berechne die Nullstelle von $f(x) = 3x - 12$.}} &
\end{gleichungsraster}
\begin{teile}
\teil Graph und Nullstelle: Zeichne den Graphen der Funktion f mit $f(x) = -4x + 6$ und berechne die Nullstelle von f. \hfill (P10 2022 OS) \\

\begin{ksys}[xmin=-4,xmax=4,ymin=-4,ymax=7] \end{ksys}
x = \leerfeld
\end{teile}
\end{aufgabe}

%% TODO Titel: Ich-kann-Satz fehlt in der Bank (Platzhalter: Kettenname) – Nr. 21
%% TODO Anweisung: ein Satz über der ersten Teilaufgabe fehlt in der Bank – Nr. 21
\begin{aufgabe}{Wertetabelle}
\begin{teile}
\teil Wertetabelle: Fülle die Tabelle zu $f(x) = -\frac{1}{2}x + 3$ aus.

\wertetabelle{x}{f(x)}{-4,-2,0,2,4}
\end{teile}
\end{aufgabe}

%% TODO Titel: Ich-kann-Satz fehlt in der Bank (Platzhalter: Kettenname) – Nr. 22
%% TODO Anweisung: ein Satz über der ersten Teilaufgabe fehlt in der Bank – Nr. 22
\begin{aufgabe}{Schnittpunkt mit y-Achse}
\begin{teile}
\teil Schnittpunkt mit der y-Achse? Gib ihn für $f(x) = 5x - 8$ an. S( \leerfeld | \leerfeld )
\end{teile}
\end{aufgabe}

%% TODO Titel: Ich-kann-Satz fehlt in der Bank (Platzhalter: Kettenname) – Nr. 23
%% TODO Anweisung: ein Satz über der ersten Teilaufgabe fehlt in der Bank – Nr. 23
\begin{aufgabe}{Funktionswert – Fehler finden, Begründen, Darstellungswechsel, Anwendung}
\begin{teile}
\teil Ich finde den Fehler: Sara soll die Nullstelle von $f(x) = 2x - 6$ angeben und schreibt: Die Nullstelle ist $-6$. Benenne den Fehler und rechne richtig.
\teil Zwei Nullstellen? Kann die Gerade $f(x) = 3x - 5$ zwei Nullstellen haben? Begründe.
\teil Gleichung aus der Tabelle: Gib die Gleichung der linearen Funktion f an.

\wertetabelle[-1,2,5,8]{x}{f(x)}{0,1,2,3}
f(x) = \leerfeld
\end{teile}
\begin{gleichungsraster}[2]
\gl{\text{Wann ist der Aufzug unten? Ein Aufzug fährt aus 36 m Höhe gleichmäßig nach unten. Seine Höhe über dem Boden (in m) nach x Sekunden ist $f(x) = -1{,}5x + 36$. Nach wie vielen Sekunden ist er unten?}} & \\
\end{gleichungsraster}
\end{aufgabe}

\clearpage
% Einheit 4 von 5 – Gleichung bestimmen (bank/lineare-funktionen/e4.jsonl, zusammenbau v0.1)
%% TODO Zweigzeile Teil 1: was hier gelernt wird (Satz in Schülersprache aus dem Einheitstitel) – Einheit 4
\einheitenkopf[e4]{Einheit 4 von 5 · Gleichung bestimmen}
\zweigzeile{ab Kl. 8 · P10}
\setcounter{aufgabe}{23}

%% TODO Titel: Ich-kann-Satz fehlt in der Bank (Platzhalter: Kettenname) – Nr. 24
%% TODO Anweisung: ein Satz über der ersten Teilaufgabe fehlt in der Bank – Nr. 24
\begin{aufgabe}{Gleichung bestimmen}
%% TODO \rechenplatz{n}: Schrittzahl des Grundfalls fehlt in der Bank (nur bei fester Schreibform, 2.3 b)
\begin{teile}
\teil Gleichung? Bestimme die Gleichung der Geraden g aus dem Graphen.

\begin{ksys}[xmin=-5,xmax=5,ymin=-5,ymax=5,ablesen] \gerade{2}{-3}{g} \end{ksys}
f(x) = \leerfeld
\end{teile}
\begin{gleichungsraster}[2]
\gl{\text{Gleichung? Eine Gerade hat die Steigung $m = 2$ und geht durch $P(4|7)$. Bestimme ihre Gleichung.}} &
\gl{\text{Gleichung aus zwei Punkten: Bestimme die Gleichung der Geraden durch $A(0|4)$ und $B(2|10)$.}} \\
\gl{\text{Gleichung aus zwei Punkten: Bestimme die Gleichung der Geraden durch $A(2|3)$ und $B(5|9)$.}} &
\gl{\text{Gleichung aus zwei Punkten: Bestimme die Gleichung der Geraden durch $A(1|6)$ und $B(3|0)$.}} \\
\end{gleichungsraster}
\begin{teile}
\teil Gerade durch A und B: Zeichne die Gerade f durch $A(-2|8)$ und $B(4|-1)$, entscheide für jede Aussage, ob sie wahr oder falsch ist, und gib eine Gleichung von f an. \hfill (P10 2025 OS) \\ \kreuz{wahr}\kreuz{falsch} f verläuft monoton steigend. \\ \kreuz{wahr}\kreuz{falsch} f schneidet die y-Achse in $(0|5)$.

\begin{ksys}[xmin=-4,xmax=5,ymin=-4,ymax=9] \end{ksys}
f(x) = \leerfeld
\end{teile}
\end{aufgabe}

%% TODO Titel: Ich-kann-Satz fehlt in der Bank (Platzhalter: Kettenname) – Nr. 25
%% TODO Anweisung: ein Satz über der ersten Teilaufgabe fehlt in der Bank – Nr. 25
\begin{aufgabe}{Gerade durch zwei Punkte zeichnen}
\begin{teile}
\teil Gerade durch zwei Punkte: Trage $A(-3|-1)$ und $B(2|4)$ ein und zeichne die Gerade durch beide Punkte.

\begin{ksys}[xmin=-4,xmax=4,ymin=-4,ymax=5] \end{ksys}
\end{teile}
\end{aufgabe}

%% TODO Titel: Ich-kann-Satz fehlt in der Bank (Platzhalter: Kettenname) – Nr. 26
%% TODO Anweisung: ein Satz über der ersten Teilaufgabe fehlt in der Bank – Nr. 26
\begin{aufgabe}{Schnittpunkt zweier Geraden rechnerisch}
\begin{gleichungsraster}[2]
\gl{\text{Schnittpunkt? Berechne den Schnittpunkt der Geraden $f(x) = 3x - 4$ und $g(x) = -x + 8$.}} & \\
\end{gleichungsraster}
\end{aufgabe}

%% TODO Titel: Ich-kann-Satz fehlt in der Bank (Platzhalter: Kettenname) – Nr. 27
%% TODO Anweisung: ein Satz über der ersten Teilaufgabe fehlt in der Bank – Nr. 27
\begin{aufgabe}{Gleichung bestimmen – Fehler finden, Begründen, Darstellungswechsel, Anwendung}
\begin{teile}
\teil Ich finde den Fehler: Für $A(2|3)$ und $B(5|12)$ rechnet Kim: \rechnung{m &= (3 - 12) : (5 - 2) \\ &= -3} Benenne den Fehler und rechne richtig.
\teil Schneiden sich die Geraden $f(x) = 3x - 7$ und $g(x) = 3x + 5$? Begründe ohne Rechnung.
\teil Gleichung zur Geraden: Bestimme die Gleichung der Geraden g aus dem Graphen.

\begin{ksys}[xmin=-5,xmax=5,ymin=-5,ymax=5,ablesen] \gerade{-0.5}{1}{g} \end{ksys}
f(x) = \leerfeld
\end{teile}
\begin{gleichungsraster}[2]
\gl{\text{Wie hoch ist die Grundgebühr? Für eine Taxifahrt von 4 km zahlt man 12{,}60 €, für 10 km 25{,}20 €. Der Preis setzt sich aus Grundgebühr und festem Preis je km zusammen.}} & \\
\end{gleichungsraster}
\end{aufgabe}

\clearpage
% Einheit 5 von 5 – Anwendungen (bank/lineare-funktionen/e5.jsonl, zusammenbau v0.1)
%% TODO Zweigzeile Teil 1: was hier gelernt wird (Satz in Schülersprache aus dem Einheitstitel) – Einheit 5
\einheitenkopf[e5]{Einheit 5 von 5 · Anwendungen}
\zweigzeile{ab Kl. 8 · P10 oft}
\setcounter{aufgabe}{27}

%% TODO Titel: Ich-kann-Satz fehlt in der Bank (Platzhalter: Kettenname) – Nr. 28
%% TODO Anweisung: ein Satz über der ersten Teilaufgabe fehlt in der Bank – Nr. 28
\begin{aufgabe}{Anwendung}
%% TODO \rechenplatz{n}: Schrittzahl des Grundfalls fehlt in der Bank (nur bei fester Schreibform, 2.3 b)
\begin{teile}
\teil Markiere: Unterstreiche den Anfangswert, kreise die Änderung je Einheit ein. Ein Handwerker berechnet 45 € für die Anfahrt und 38 € je Arbeitsstunde.
\teil Welche Gleichung passt? Ein Fahrradverleih nimmt 8 € Grundgebühr und 3 € je Stunde. y ist der Preis in €, x die Anzahl der Einheiten. Kreuze an. \\ \kreuz{$y = 8x + 3$} \\ \kreuz{$y = 3x + 8$} \\ \kreuz{$y = 11x$} \\ \kreuz{$y = 3x - 8$}
\teil Stelle eine Gleichung auf: Ein Kino-Abo kostet einmalig 10 € und 6 € je Film. y sind die Kosten in €, x die Anzahl der Filme. y = \leerfeld
\teil Wie viel nach 12 Wochen? Tom hat 85 € gespart. Jede Woche kommen 15 € dazu. \leerfeld[€]
\teil Welcher Tarif ist für 60 km günstiger? Tarif A: 20 € Grundgebühr und 0{,}30 € je km. Tarif B: 0{,}70 € je km ohne Grundgebühr.
\end{teile}
\begin{gleichungsraster}[2]
\gl{\text{Welches Mietauto? Angebot 1: 32 € Miete pro Tag, 300 km insgesamt frei, jeder weitere km 0{,}30 €. Angebot 2: einmalig 95 € für eine Woche, jeder km 0{,}12 €. Jana mietet 4 Tage und fährt 420 km. Berechne die Gesamtkosten beider Angebote, entscheide, welches günstiger ist, und stelle für Angebot 2 eine Gleichung für die Kosten y (in €) bei x km für eine Woche auf. (P10 2023 OS)}} & \\
\end{gleichungsraster}
\end{aufgabe}

%% TODO Titel: Ich-kann-Satz fehlt in der Bank (Platzhalter: Kettenname) – Nr. 29
%% TODO Anweisung: ein Satz über der ersten Teilaufgabe fehlt in der Bank – Nr. 29
\begin{aufgabe}{Graph zu Tarif zuordnen}
\begin{teile}
\teil Welcher Graph passt zum Tarif? Tarif: 10 € Grundgebühr und 2 € je Stunde. Gib den passenden Graphen an.

\begin{ksys}[xmin=0,xmax=8,ymin=0,ymax=40,xstep=1,ystep=5,ablesen] \gerade{2}{0}{A} \gerade{2}{10}{B} \gerade{10}{2}{C} \end{ksys}
Graph \leerfeld
\end{teile}
\end{aufgabe}

%% TODO Titel: Ich-kann-Satz fehlt in der Bank (Platzhalter: Kettenname) – Nr. 30
%% TODO Anweisung: ein Satz über der ersten Teilaufgabe fehlt in der Bank – Nr. 30
\begin{aufgabe}{Situation zu Gleichung beschreiben}
\begin{teile}
\teil Welche Situation passt? Beschreibe eine Situation zur Gleichung $y = 2{,}5x + 4$ und sage, was x und y bedeuten.
\end{teile}
\end{aufgabe}

%% TODO Titel: Ich-kann-Satz fehlt in der Bank (Platzhalter: Kettenname) – Nr. 31
%% TODO Anweisung: ein Satz über der ersten Teilaufgabe fehlt in der Bank – Nr. 31
\begin{aufgabe}{Anwendung – Fehler finden, Begründen, Darstellungswechsel, Anwendung}
\begin{teile}
\teil Ich finde den Fehler: Ein Verleih kostet 6 € Grundgebühr und 2{,}50 € je Stunde. Max schreibt $y = 6x + 2{,}5$. Benenne den Fehler und gib die richtige Gleichung an.
\teil Welcher Tarif ist immer günstiger? Tarif A: 10 € Grundgebühr und 0{,}20 € je Minute. Tarif B: 10 € Grundgebühr und 0{,}15 € je Minute. Begründe ohne Rechnung.
\teil Gleichung aus der Tabelle: Die Tabelle zeigt die Kosten y (in €) eines Handytarifs für x GB. Gib die Gleichung des Tarifs an.

\wertetabelle[12,15,18,21]{x}{y}{0,2,4,6}
y = \leerfeld
\end{teile}
\begin{gleichungsraster}[2]
\gl{\text{Wie weit reicht das Geld? Für die Klassenfahrt kostet ein Bus 180 € Grundpreis und 1{,}90 € je km. Die Klasse hat 600 € für den Bus. Wie viele Kilometer (hin und zurück) sind höchstens drin?}} & \\
\end{gleichungsraster}
\end{aufgabe}

% Abhakseite „Das kann ich“ – Zone nur im Gesamt (\mitzone)
%% TODO Ich-kann-Sätze: Platzhalter wie in den Aufgabendateien
\begin{abhakseite}
\ifdefined\mitzone
\abhakgruppe{Kennst du schon}
\abhak{1}{Koordinaten lesen und eintragen, 4 Quadranten, (x|y)-Reihenfolge}
\abhak{2}{Proportionale Zuordnung erkennen und hochrechnen (Dreisatz)}
\abhak{3}{Negative Zahlen multiplizieren und dividieren}
\abhak{4}{Brüche als Steigung (½, −¾), Bruch mal ganze Zahl}
\abhak{5}{Lineare Gleichung zweischrittig lösen (null gleich einem zweischrittigen Term)}
\abhak{6}{Terme zusammenfassen}
\abhak{7}{Negative Zahlen multiplizieren und dividieren – Fehler finden}
\abhak{8}{Negative Zahlen multiplizieren und dividieren – selbst rechnen}
\fi
\abhakgruppe{Einheit 1 · Proportionale Funktion}
\abhak{9}{Proportionale Funktion}
\abhak{10}{Dreisatz als Kontrolle}
\abhak{11}{Proportionale Funktion – Fehler finden, Begründen, Darstellungswechsel, Anwendung}
\abhakgruppe{Einheit 2 · Lineare Funktion f(x) = m·x + n}
\abhak{12}{m und n in der Gleichung markieren}
\abhak{13}{Steigungsrichtung ohne Zeichnen}
\abhak{14}{Steigungsdreieck lesen}
\abhak{15}{Graph zeichnen}
\abhak{16}{Ablesen}
\abhak{17}{Parameter deuten (steigend, fallend, parallel, Sonderfall m = 0)}
\abhak{18}{Graph zeichnen – Fehler finden, Begründen, Darstellungswechsel}
\abhakgruppe{Einheit 3 · Punkte und Werte}
\abhak{19}{Punkt in Gleichung einsetzen}
\abhak{20}{Funktionswert}
\abhak{21}{Wertetabelle}
\abhak{22}{Schnittpunkt mit y-Achse}
\abhak{23}{Funktionswert – Fehler finden, Begründen, Darstellungswechsel, Anwendung}
\abhakgruppe{Einheit 4 · Gleichung bestimmen}
\abhak{24}{Gleichung bestimmen}
\abhak{25}{Gerade durch zwei Punkte zeichnen}
\abhak{26}{Schnittpunkt zweier Geraden rechnerisch}
\abhak{27}{Gleichung bestimmen – Fehler finden, Begründen, Darstellungswechsel, Anwendung}
\abhakgruppe{Einheit 5 · Anwendungen}
\abhak{28}{Anwendung}
\abhak{29}{Graph zu Tarif zuordnen}
\abhak{30}{Situation zu Gleichung beschreiben}
\abhak{31}{Anwendung – Fehler finden, Begründen, Darstellungswechsel, Anwendung}
\end{abhakseite}

\end{document}
